% NGÂN HÀNG TUYỂN SINH 10 – BÀI 6

\begin{bt}%[TS10 6.3]%[ma:e550352a]
% nguon: Đề thi chính thức Tuyển sinh 10 TP.HCM năm học 2025-2026 – Bài 6
Từ vị trí $A$ của một công viên có dạng hình vuông $ABCD$ cạnh $a\text{ (km)}$, hai bạn Hoà và Bình bắt đầu chạy bộ cùng lúc với vận tốc không đổi dọc theo các cạnh của hình vuông và theo hai hướng khác nhau (Hoà chạy theo hướng $A \to B \to C \to D$, Bình chạy theo hướng $A \to D \to C \to B$). Biết rằng hai bạn gặp nhau lần thứ nhất tại vị trí $E$ trên cạnh $CD$ cách $A$ một khoảng bằng $1\text{ km}$ ($AE = 1\text{ km}$) và gặp lại nhau lần thứ hai tại vị trí $F$ trên cạnh $AB$ cách $A$ một khoảng bằng $0{,}4\text{ km}$ ($AF = 0{,}4\text{ km}$). Gọi $x, y\text{ (km/h)}$ lần lượt là vận tốc của Hoà và Bình.
\begin{enumerate}[a)]
	\item Chứng minh rằng $\dfrac{x}{y} = \dfrac{AB + BC + CE}{AD + DE}$.
	\item Tìm giá trị của $a$.
\end{enumerate}
\loigiai{
\begin{enumerate}[a)]
	\item Hai bạn xuất phát cùng lúc từ $A$. Khi gặp nhau lần thứ nhất tại $E$ trên cạnh $CD$:
	\begin{itemize}
		\item Quãng đường Hoà đã chạy là: $s_1 = AB + BC + CE$.
		\item Quãng đường Bình đã chạy là: $s_2 = AD + DE$.
	\end{itemize}
	Vì thời gian chạy của hai bạn từ lúc xuất phát đến lúc gặp nhau lần đầu là như nhau nên:
	$t_1 = \dfrac{s_1}{x} = \dfrac{s_2}{y} \Leftrightarrow \dfrac{AB + BC + CE}{x} = \dfrac{AD + DE}{y}$.\\
	Suy ra $\dfrac{x}{y} = \dfrac{AB + BC + CE}{AD + DE}$.
	\item Chu vi công viên hình vuông là $4a\text{ (km)}$.\\
	Tổng quãng đường hai bạn chạy được tính đến lần gặp thứ nhất là:
	$s_1 + s_2 = (AB + BC + CE) + (AD + DE) = AB + BC + CD + DA = 4a\text{ (km)}$.\\
	Do đó thời gian chạy đến lần gặp thứ nhất là:
	$t_1 = \dfrac{4a}{x + y}$.\\
	Từ lần gặp thứ nhất đến lần gặp thứ hai tại $F$, hai bạn tiếp tục chạy ngược chiều quanh chu vi hình vuông và gặp nhau, nên tổng quãng đường hai bạn chạy thêm đúng bằng một vòng chu vi hình vuông là $4a\text{ (km)}$.\\
	Do đó, thời gian từ lúc xuất phát đến khi gặp nhau lần thứ hai là:
	$t_2 = \dfrac{4a + 4a}{x + y} = \dfrac{8a}{x + y} = 2t_1$.\\
	Vì vận tốc của Bình là không đổi nên quãng đường Bình chạy đến lần gặp thứ hai gấp đôi quãng đường Bình chạy đến lần gặp thứ nhất.\\
	Tại lần gặp thứ hai ở $F$ trên cạnh $AB$ cách $A$ là $AF = 0{,}4\text{ km}$, quãng đường Bình đã chạy từ $A$ qua $D, C, B$ đến $F$ là:
	$s'_2 = AD + DC + CB + BF = 4a - AF = 4a - 0{,}4\text{ (km)}$.\\
	Do $s'_2 = 2s_2$, ta có:
	$4a - 0{,}4 = 2(AD + DE) = 2(a + DE) \Leftrightarrow 4a - 0{,}4 = 2a + 2DE \Leftrightarrow 2DE = 2a - 0{,}4 \Leftrightarrow DE = a - 0{,}2$.\\
	Tam giác $ADE$ vuông tại $D$ có $AD = a$, $DE = a - 0{,}2$ và cạnh huyền $AE = 1\text{ km}$. Áp dụng định lí Pythagore:
	$AD^2 + DE^2 = AE^2 \Leftrightarrow a^2 + (a - 0{,}2)^2 = 1^2 \Leftrightarrow a^2 + a^2 - 0{,}4a + 0{,}04 = 1 \Leftrightarrow 2a^2 - 0{,}4a - 0{,}96 = 0$.\\
	Chia cả hai vế cho $2$: $a^2 - 0{,}2a - 0{,}48 = 0$.\\
	Ta có $\Delta' = (-0{,}1)^2 - 1\cdot(-0{,}48) = 0{,}01 + 0{,}48 = 0{,}49 = (0{,}7)^2$.\\
	Phương trình có hai nghiệm:
	$a_1 = 0{,}1 + 0{,}7 = 0{,}8$ (thoả mãn $a > 0{,}2$),\\
	$a_2 = 0{,}1 - 0{,}7 = -0{,}6$ (loại vì $a > 0$).\\
	Vậy $a = 0{,}8\text{ (km)}$.
\end{enumerate}
}
\end{bt}

\begin{bt}%[TS10 6.2]%[ma:17ea06ce]
% nguon: Đề tham khảo Tuyển sinh 10 TP.HCM 2025-2026 – Quận 3 – Đề 2 – Bài 6
Nhân dịp ngày Nhà giáo Việt Nam, nhà trường tổ chức cuộc thi kéo co. Giáo viên chủ nhiệm lớp 9B đã chọn ra $\dfrac{1}{4}$ số học sinh nam và $\dfrac{1}{3}$ số học sinh nữ của lớp tham gia kéo co, biết rằng số học sinh nam và số học sinh nữ tham gia kéo co là bằng nhau. Toàn bộ $20$ học sinh còn lại của lớp tham gia đội cổ động. Hãy tìm số học sinh nam và số học sinh nữ của lớp 9B.
\loigiai{
Gọi số học sinh nam và số học sinh nữ của lớp 9B lần lượt là $x$ và $y$ ($x, y \in \mathbb{N}^*$, $x$ chia hết cho $4$ và $y$ chia hết cho $3$).\\
Số học sinh nam tham gia kéo co là $\dfrac{1}{4}x$ (học sinh).\\
Số học sinh nữ tham gia kéo co là $\dfrac{1}{3}y$ (học sinh).\\
Vì số học sinh nam và số học sinh nữ tham gia kéo co bằng nhau nên ta có phương trình:
$\dfrac{1}{4}x = \dfrac{1}{3}y \Leftrightarrow 3x - 4y = 0$.\\
Số học sinh nam làm cổ động viên là: $x - \dfrac{1}{4}x = \dfrac{3}{4}x$ (học sinh).\\
Số học sinh nữ làm cổ động viên là: $y - \dfrac{1}{3}y = \dfrac{2}{3}y$ (học sinh).\\
Vì đội cổ động có tổng cộng $20$ học sinh nên ta có phương trình:
$\dfrac{3}{4}x + \dfrac{2}{3}y = 20 \Leftrightarrow 9x + 8y = 240$.\\
Ta có hệ phương trình:
$\begin{cases} 3x - 4y = 0 \\ 9x + 8y = 240 \end{cases} \Leftrightarrow \begin{cases} 6x - 8y = 0 \\ 9x + 8y = 240 \end{cases} \Leftrightarrow \begin{cases} 15x = 240 \\ y = \dfrac{3}{4}x \end{cases} \Leftrightarrow \begin{cases} x = 16 \\ y = 12 \end{cases}$ (thoả mãn điều kiện).\\
Vậy lớp 9B có $16$ học sinh nam và $12$ học sinh nữ.
}
\end{bt}

\begin{bt}%[TS10 6.7]%[ma:439a2838]
% nguon: Đề thi tham khảo Tuyển sinh 10 TP.HCM – Phụ lục Sở GD&ĐT TP.HCM – Bài 6
Thép không gỉ Ferritic là họ thép hợp kim có chứa từ $12\%$ đến $27\%$ crôm. Một nhà máy luyện thép hiện có sẵn một lượng hợp kim thép chứa $10\%$ crôm và một lượng hợp kim thép chứa $30\%$ crôm. Giả sử trong quá trình luyện thép các nguyên liệu không bị hao hụt.
\begin{enumerate}[a)]
	\item Tính khối lượng hợp kim thép mỗi loại từ hai loại thép trên dùng để luyện được $500\text{ tấn}$ thép chứa $16\%$ crôm.
	\item Nhà máy dự định luyện ra loại thép không gỉ Ferritic từ $100\text{ tấn}$ thép chứa $10\%$ crôm và $x\text{ tấn}$ thép chứa $30\%$ crôm. Hỏi $x$ nằm trong khoảng nào?
\end{enumerate}
\loigiai{
\begin{enumerate}[a)]
	\item Gọi khối lượng hợp kim thép chứa $10\%$ crôm và chứa $30\%$ crôm cần dùng lần lượt là $u$ và $v$ (tấn, điều kiện: $0 < u, v < 500$).\\
	Theo đề bài, tổng khối lượng thép luyện được là $500\text{ tấn}$, ta có:
	$u + v = 500$.\\
	Lượng crôm nguyên chất có trong sản phẩm là $500 \cdot 16\% = 80\text{ (tấn)}$, ta có phương trình:
	$10\% u + 30\% v = 80 \Leftrightarrow 0{,}1u + 0{,}3v = 80 \Leftrightarrow u + 3v = 800$.\\
	Ta có hệ phương trình:
	$\begin{cases} u + v = 500 \\ u + 3v = 800 \end{cases} \Leftrightarrow \begin{cases} 2v = 300 \\ u = 500 - v \end{cases} \Leftrightarrow \begin{cases} v = 150 \\ u = 350 \end{cases}$ (thoả mãn điều kiện).\\
	Vậy cần dùng $350\text{ tấn}$ thép chứa $10\%$ crôm và $150\text{ tấn}$ thép chứa $30\%$ crôm.
	\item Khi trộn $100\text{ tấn}$ thép chứa $10\%$ crôm với $x\text{ tấn}$ thép chứa $30\%$ crôm ($x > 0$):
	\begin{itemize}
		\item Tổng khối lượng thép thu được là: $100 + x\text{ (tấn)}$.
		\item Tổng lượng crôm nguyên chất thu được là: $100\cdot 10\% + x\cdot 30\% = 10 + 0{,}3x\text{ (tấn)}$.
	\end{itemize}
	Tỉ lệ phần trăm crôm trong hợp kim mới là: $\dfrac{10 + 0{,}3x}{100 + x}$.\\
	Để hợp kim thu được là thép không gỉ Ferritic (chứa từ $12\%$ đến $27\%$ crôm), ta có bất phương trình kép:
	$12\% \le \dfrac{10 + 0{,}3x}{100 + x} \le 27\% \Leftrightarrow 0{,}12 \le \dfrac{10 + 0{,}3x}{100 + x} \le 0{,}27$.\\
	Vì $x > 0$ nên $100 + x > 0$, nhân cả hai vế với $(100 + x)$:
	\begin{itemize}
		\item $0{,}12(100 + x) \le 10 + 0{,}3x \Leftrightarrow 12 + 0{,}12x \le 10 + 0{,}3x \Leftrightarrow 0{,}18x \ge 2 \Leftrightarrow x \ge \dfrac{2}{0{,}18} = \dfrac{100}{9} \approx 11{,}11\text{ (tấn)}$.
		\item $10 + 0{,}3x \le 0{,}27(100 + x) \Leftrightarrow 10 + 0{,}3x \le 27 + 0{,}27x \Leftrightarrow 0{,}03x \le 17 \Leftrightarrow x \le \dfrac{17}{0{,}03} = \dfrac{1700}{3} \approx 566{,}67\text{ (tấn)}$.
	\end{itemize}
	Vậy khối lượng $x$ cần dùng nằm trong khoảng $\dfrac{100}{9} \le x \le \dfrac{1700}{3}$ (tấn), tức là khoảng từ $11{,}11\text{ tấn}$ đến $566{,}67\text{ tấn}$.
\end{enumerate}
}
\end{bt}

\begin{bt}%[TS10 6.2]%[ma:d63ce6d2]
% nguon: Đề tham khảo Tuyển sinh 10 TP.HCM 2025-2026 – Quận 1 – Đề 1 – Bài 6
Anh Bình đang có kế hoạch tiết kiệm để mua một chiếc xe ô tô cũ có giá $250\text{ triệu đồng}$ trong vòng một năm ($12\text{ tháng}$). Hiện tại, anh Bình đã có $50\text{ triệu đồng}$ trong tài khoản tiết kiệm.
\begin{enumerate}[a)]
	\item Hãy xác định số tiền tối thiểu anh Bình cần tiết kiệm mỗi tháng để đạt được mục tiêu mua xe sau $12\text{ tháng}$.
	\item Nếu vào tháng thứ $7$, anh Bình phải dùng $15\text{ triệu đồng}$ từ khoản tiết kiệm để gửi về hỗ trợ gia đình, thì trung bình mỗi tháng trong $12\text{ tháng}$ anh Bình cần tiết kiệm ít nhất bao nhiêu tiền để vẫn đạt được mục tiêu mua xe?
\end{enumerate}
\loigiai{
\begin{enumerate}[a)]
	\item Số tiền anh Bình còn thiếu để đủ $250\text{ triệu đồng}$ mua xe là:
	$250 - 50 = 200\text{ (triệu đồng)}$.\\
	Số tiền tối thiểu anh Bình cần tiết kiệm mỗi tháng trong $12\text{ tháng}$ là:
	$\dfrac{200}{12} = \dfrac{50}{3}\text{ (triệu đồng)} \approx 16{,}67\text{ triệu đồng} = 16\,666\,667\text{ đồng}$.
	\item Khi rút $15\text{ triệu đồng}$ gửi về cho gia đình, tổng số tiền anh Bình cần tiết kiệm trong cả năm là:
	$200 + 15 = 215\text{ (triệu đồng)}$.\\
	Số tiền ít nhất anh Bình cần tiết kiệm trung bình mỗi tháng là:
	$\dfrac{215}{12}\text{ (triệu đồng)} \approx 17{,}92\text{ triệu đồng} = 17\,916\,667\text{ đồng}$.
\end{enumerate}
}
\end{bt}

\begin{bt}%[TS10 6.1]%[ma:5469bf96]
% nguon: Đề tham khảo Tuyển sinh 10 TP.HCM 2025-2026 – Quận 1 – Đề 2 – Bài 6
Chị Ba muốn mở một cửa hàng bán trà sữa với giá niêm yết là $24\,000\text{ đồng/ly}$. Nhân dịp khai trương, chị Ba cân nhắc hai hình thức khuyến mãi:
\begin{itemize}
	\item Phương án 1: Giảm giá $30\%$ cho mỗi ly trà sữa.
	\item Phương án 2: Mua $2$ ly được tặng thêm $1$ ly cùng loại.
\end{itemize}
Theo em, nếu hiệu ứng thu hút khách hàng của hai phương án là như nhau thì chị Ba nên chọn phương án nào để có lợi hơn cho cửa hàng? Vì sao?
\loigiai{
Xét số tiền thu được của cửa hàng đối với từng phương án trên mỗi đơn vị sản phẩm bán ra:
\begin{itemize}
	\item Theo Phương án 1 (giảm giá $30\%$ mỗi ly):\\
	Giá bán thực tế mỗi ly trà sữa là:
	$24\,000 \cdot (100\% - 30\%) = 24\,000 \cdot 0{,}7 = 16\,800\text{ (đồng/ly)}$.\\
	Do đó, mỗi ly bán ra cửa hàng thu về $16\,800\text{ đồng}$.
	\item Theo Phương án 2 (mua $2$ tặng $1$):\\
	Khi khách hàng nhận $3$ ly trà sữa thì chỉ phải trả tiền $2$ ly. Số tiền khách trả là:
	$2 \cdot 24\,000 = 48\,000\text{ (đồng)}$.\\
	Giá bán trung bình mỗi ly trà sữa mà cửa hàng giao cho khách là:
	$\dfrac{48\,000}{3} = 16\,000\text{ (đồng/ly)}$.\\
	Do đó, mỗi ly xuất ra cửa hàng chỉ thu về $16\,000\text{ đồng}$.
\end{itemize}
Vì $16\,800\text{ đồng} > 16\,000\text{ đồng}$ (chênh lệch $800\text{ đồng}$ trên mỗi ly) nên Phương án 1 mang lại doanh thu trên từng ly cao hơn so với Phương án 2.\\
Vậy chị Ba nên chọn Phương án 1 để có lợi hơn cho cửa hàng.
}
\end{bt}

\begin{bt}%[TS10 6.1]%[ma:cd90e865]
% nguon: Đề tham khảo Tuyển sinh 10 TP.HCM 2025-2026 – Quận 4 – Đề 1 – Bài 6
Cô Mai nhập $210\text{ bao gạo}$ gồm hai loại: bao gạo loại 1 (mỗi bao nặng $25\text{ kg}$, bán với giá $22\,500\text{ đồng/kg}$) và bao gạo loại 2 (mỗi bao nặng $50\text{ kg}$, bán với giá $18\,200\text{ đồng/kg}$) với tổng chi phí nhập hàng là $99\,700\,000\text{ đồng}$. Do thời tiết ẩm ướt, $5\%$ khối lượng gạo loại 1 và $6\%$ khối lượng gạo loại 2 bị hỏng không thể bán được. Tổng khối lượng gạo còn lại bán được là $6{,}15\text{ tấn} = 6150\text{ kg}$. Hỏi sau khi bán hết toàn bộ số gạo còn lại, cô Mai lãi hay lỗ bao nhiêu tiền?
\loigiai{
Gọi số bao gạo loại 1 và loại 2 cô Mai đã nhập lần lượt là $x$ và $y$ ($x, y \in \mathbb{N}^*, x, y < 210$).\\
Vì tổng số bao gạo nhập về là $210$ bao nên ta có phương trình:
$x + y = 210$.\\
Tổng khối lượng gạo ban đầu của loại 1 là $25x\text{ (kg)}$, sau khi hỏng $5\%$ thì khối lượng bán được là:
$25x \cdot (100\% - 5\%) = 25x \cdot 0{,}95 = 23{,}75x\text{ (kg)}$.\\
Tổng khối lượng gạo ban đầu của loại 2 là $50y\text{ (kg)}$, sau khi hỏng $6\%$ thì khối lượng bán được là:
$50y \cdot (100\% - 6\%) = 50y \cdot 0{,}94 = 47y\text{ (kg)}$.\\
Vì tổng khối lượng gạo còn lại bán được là $6150\text{ kg}$ nên ta có phương trình:
$23{,}75x + 47y = 6150$.\\
Ta có hệ phương trình:
$\begin{cases} x + y = 210 \\ 23{,}75x + 47y = 6150 \end{cases} \Leftrightarrow \begin{cases} y = 210 - x \\ 23{,}75x + 47(210 - x) = 6150 \end{cases} \Leftrightarrow \begin{cases} 23{,}25x = 3720 \\ y = 210 - x \end{cases} \Leftrightarrow \begin{cases} x = 160 \\ y = 50 \end{cases}$ (thoả mãn điều kiện).\\
Khối lượng gạo loại 1 bán được là: $160 \cdot 23{,}75 = 3800\text{ (kg)}$.\\
Doanh thu từ gạo loại 1 là: $3800 \cdot 22\,500 = 85\,500\,000\text{ (đồng)}$.\\
Khối lượng gạo loại 2 bán được là: $50 \cdot 47 = 2350\text{ (kg)}$.\\
Doanh thu từ gạo loại 2 là: $2350 \cdot 18\,200 = 42\,770\,000\text{ (đồng)}$.\\
Tổng doanh thu thu được từ việc bán gạo là:
$85\,500\,000 + 42\,770\,000 = 128\,270\,000\text{ (đồng)}$.\\
Số tiền lãi của cô Mai sau khi trừ chi phí nhập hàng là:
$128\,270\,000 - 99\,700\,000 = 28\,570\,000\text{ (đồng)}$.\\
Vậy cô Mai có lãi và số tiền lãi là $28\,570\,000\text{ đồng}$.
}
\end{bt}

\begin{bt}%[TS10 6.8]%[ma:0f1866c1]
% nguon: Đề tham khảo Tuyển sinh 10 TP.HCM 2025-2026 – Quận 4 – Đề 2 – Bài 6
Điểm trung bình môn Toán của $100$ học sinh trong hai lớp 9A và 9B là $7{,}608$. Hãy tính điểm trung bình của học sinh mỗi lớp, biết rằng số học sinh của lớp 9A nhiều hơn lớp 9B là $2$ học sinh và điểm trung bình của học sinh lớp 9B bằng $\dfrac{9}{10}$ điểm trung bình của học sinh lớp 9A.
\loigiai{
Gọi số học sinh của lớp 9A và lớp 9B lần lượt là $x$ và $y$ ($x, y \in \mathbb{N}^*, x, y < 100$).\\
Vì tổng số học sinh hai lớp là $100$ và lớp 9A nhiều hơn lớp 9B là $2$ học sinh, nên ta có hệ phương trình:
$\begin{cases} x + y = 100 \\ x - y = 2 \end{cases} \Leftrightarrow \begin{cases} x = 51 \\ y = 49. \end{cases}$\\
Gọi điểm trung bình của học sinh lớp 9A là $u$ ($u > 0$).\\
Khi đó điểm trung bình của học sinh lớp 9B là: $v = \dfrac{9}{10}u = 0{,}9u$.\\
Tổng điểm thi môn Toán của cả hai lớp là $100 \cdot 7{,}608 = 760{,}8\text{ (điểm)}$.\\
Ta có phương trình:
$51u + 49(0{,}9u) = 760{,}8 \Leftrightarrow 51u + 44{,}1u = 760{,}8 \Leftrightarrow 95{,}1u = 760{,}8 \Leftrightarrow u = \dfrac{760{,}8}{95{,}1} = 8$.\\
Suy ra điểm trung bình của lớp 9B là: $v = 0{,}9 \cdot 8 = 7{,}2$.\\
Vậy điểm trung bình của lớp 9A là $8{,}0$ điểm và lớp 9B là $7{,}2$ điểm.
}
\end{bt}

\begin{bt}%[TS10 6.2]%[ma:415f5b68]
% nguon: Đề tham khảo Tuyển sinh 10 TP.HCM 2025-2026 – Quận 5 – Đề 1 – Bài 6
Nhân dịp tết Nguyên đán, một trường trung học cơ sở đã tổ chức cho học sinh các khối lớp 6, 7, 8 và 9 quyên góp tiền ủng hộ tự nguyện nhằm góp phần nhỏ giúp đỡ những học sinh có hoàn cảnh khó khăn. Kết thúc đợt quyên góp, bình quân mỗi học sinh của khối lớp 9 quyên góp được $22\,000$ đồng; bình quân mỗi học sinh của ba khối lớp 6, 7 và 8 quyên được $18\,000$ đồng. Hỏi bình quân mỗi học sinh của trường quyên góp được bao nhiêu tiền? Biết tổng số học sinh của ba khối lớp 6, 7 và 8 gấp 3 lần tổng số của học sinh khối lớp 9.
\loigiai{
Gọi $x$ là tổng số học sinh khối lớp 9 ($x \in \mathbb{N}^*$).\\
Vì tổng số học sinh của ba khối 6, 7 và 8 gấp 3 lần số học sinh khối 9 nên số học sinh ba khối 6, 7 và 8 là $3x$ (học sinh).\\
Tổng số học sinh toàn trường là: $x + 3x = 4x$ (học sinh).\\
Tổng số tiền học sinh khối lớp 9 quyên góp được là:
$22\,000 \cdot x$ (đồng).\\
Tổng số tiền học sinh ba khối 6, 7 và 8 quyên góp được là:
$18\,000 \cdot 3x = 54\,000x$ (đồng).\\
Tổng số tiền toàn trường quyên góp được là:
$22\,000x + 54\,000x = 76\,000x$ (đồng).\\
Bình quân mỗi học sinh toàn trường quyên góp được số tiền là:
$\dfrac{76\,000x}{4x} = 19\,000$ (đồng).\\
Vậy bình quân mỗi học sinh của trường quyên góp được $19\,000$ đồng.
}
\end{bt}

\begin{bt}%[TS10 6.1]%[ma:812480b6]
% nguon: Đề tham khảo Tuyển sinh 10 TP.HCM 2025-2026 – Quận 5 – Đề 2 – Bài 6
Ông Năm dự định cải tạo một mảnh vườn hình chữ nhật có chiều dài bằng $2{,}5$ lần chiều rộng. Ông nhận thấy rằng, nếu đào một cái hồ có mặt hồ là hình chữ nhật thì sẽ chiếm mất $3\%$ diện tích mảnh vườn. Còn nếu giảm chiều dài mặt hồ $5\text{ m}$ và tăng chiều rộng mặt hồ $2\text{ m}$ thì mặt hồ là hình vuông và diện tích mặt hồ giảm được $20\text{ m}^2$. Tính chiều dài, chiều rộng của mảnh vườn hình chữ nhật.
\loigiai{
Gọi chiều rộng của mặt hồ ban đầu là $x\text{ (m)}$ và chiều dài của mặt hồ ban đầu là $y\text{ (m)}$ ($y > 5, x > 0, y > x$).\\
Khi giảm chiều dài $5\text{ m}$ và tăng chiều rộng $2\text{ m}$ thì mặt hồ trở thành hình vuông, ta có phương trình:
$y - 5 = x + 2 \Leftrightarrow y - x = 7 \Leftrightarrow y = x + 7$.\\
Diện tích mặt hồ ban đầu là $x \cdot y\text{ (m}^2\text{)}$.\\
Diện tích mặt hồ sau khi thay đổi kích thước là $(y - 5)(x + 2) = (x + 2)^2\text{ (m}^2\text{)}$.\\
Vì diện tích mặt hồ mới giảm $20\text{ m}^2$ so với ban đầu nên ta có phương trình:
$\begin{aligned}
xy - (x + 2)^2 = 20 &\Leftrightarrow x(x + 7) - (x^2 + 4x + 4) = 20 \\
&\Leftrightarrow x^2 + 7x - x^2 - 4x - 4 = 20 \\
&\Leftrightarrow 3x = 24 \Leftrightarrow x = 8\text{ (thỏa mãn)}.
\end{aligned}$\\
Suy ra chiều rộng của mặt hồ là $8\text{ m}$ và chiều dài của mặt hồ là $y = 8 + 7 = 15\text{ m}$.\\
Diện tích của mặt hồ là:
$S_{\text{hồ}} = 8 \cdot 15 = 120\text{ (m}^2\text{)}$.\\
Vì diện tích mặt hồ chiếm $3\%$ diện tích mảnh vườn nên diện tích mảnh vườn là:
$S_{\text{vườn}} = \dfrac{120}{3\%} = \dfrac{120}{0{,}03} = 4\,000\text{ (m}^2\text{)}$.\\
Gọi chiều rộng của mảnh vườn là $W\text{ (m)}$ ($W > 0$). Chiều dài của mảnh vườn là $2{,}5W\text{ (m)}$.\\
Ta có phương trình diện tích mảnh vườn:
$W \cdot 2{,}5W = 4\,000 \Leftrightarrow 2{,}5W^2 = 4\,000 \Leftrightarrow W^2 = 1\,600 \Rightarrow W = 40\text{ (m)}$.\\
Chiều dài của mảnh vườn là:
$2{,}5 \cdot 40 = 100\text{ (m)}$.\\
Vậy mảnh vườn có chiều rộng là $40\text{ m}$ và chiều dài là $100\text{ m}$.
}
\end{bt}

\begin{bt}%[TS10 6.8]%[ma:cdfd8eb2]
% nguon: Đề tham khảo Tuyển sinh 10 TP.HCM 2025-2026 – Quận 5 – Đề 3 – Bài 6
Một bài thi trắc nghiệm gồm 20 câu hỏi, nếu trả lời đúng thì được 5 điểm, nếu trả lời sai thì bị trừ 2 điểm, nếu bỏ qua câu trả lời thì được 0 điểm.
\begin{enumerate}[a)]
    \item Bạn Dũng tham dự bài thi và được 47 điểm. Biết Dũng bỏ qua 5 câu. Hỏi Dũng trả lời đúng mấy câu và trả lời sai mấy câu?
    \item Bạn Hoa tham dự bài thi và được 59 điểm. Hỏi Hoa trả lời đúng mấy câu, trả lời sai mấy câu và bỏ qua bao nhiêu câu?
\end{enumerate}
\loigiai{
\begin{enumerate}[a)]
    \item Gọi $x, y$ lần lượt là số câu Dũng trả lời đúng và trả lời sai ($x, y \in \mathbb{N}$).\\
    Vì Dũng bỏ qua 5 câu nên tổng số câu Dũng đã làm là:
    $x + y = 20 - 5 = 15$.\\
    Số điểm Dũng đạt được là 47 điểm nên ta có phương trình:
    $5x - 2y = 47$.\\
    Ta có hệ phương trình:
    $\left\{\begin{aligned} &x + y = 15 \\ &5x - 2y = 47 \end{aligned}\right. \Leftrightarrow \left\{\begin{aligned} &2x + 2y = 30 \\ &5x - 2y = 47 \end{aligned}\right. \Leftrightarrow \left\{\begin{aligned} &7x = 77 \\ &y = 15 - x \end{aligned}\right. \Leftrightarrow \left\{\begin{aligned} &x = 11 \\ &y = 4. \end{aligned}\right.$\\
    Vậy bạn Dũng trả lời đúng 11 câu và trả lời sai 4 câu.
    \item Gọi $a, b, c$ lần lượt là số câu Hoa trả lời đúng, sai và bỏ qua ($a, b, c \in \mathbb{N}, a + b + c = 20$).\\
    Theo bài ra ta có phương trình điểm số:
    $5a - 2b = 59 \Leftrightarrow 5a = 59 + 2b$.\\
    Vì $b \ge 0$ nên $5a \ge 59 \Rightarrow a \ge 12$.\\
    Mà tổng số câu hỏi là 20 nên $a \le 20$.\\
    Mặt khác, $59 + 2b$ là số lẻ nên $5a$ phải là số lẻ, suy ra $a$ là số lẻ.\\
    Do đó $a \in \{13; 15; 17; 19\}$.
    \begin{itemize}
        \item Với $a = 13 \Rightarrow 2b = 5 \cdot 13 - 59 = 65 - 59 = 6 \Rightarrow b = 3$.\\
        Khi đó số câu bỏ qua là $c = 20 - (13 + 3) = 4$ (thỏa mãn).
        \item Với $a = 15 \Rightarrow 2b = 5 \cdot 15 - 59 = 75 - 59 = 16 \Rightarrow b = 8$.\\
        Khi đó $c = 20 - (15 + 8) = -3 < 0$ (loại).
        \item Với $a \ge 17$ thì $b \ge \dfrac{5 \cdot 17 - 59}{2} = 13 \Rightarrow a + b \ge 30 > 20$ (loại).
    \end{itemize}
    Vậy bạn Hoa trả lời đúng 13 câu, trả lời sai 3 câu và bỏ qua 4 câu.
\end{enumerate}
}
\end{bt}

\begin{bt}%[TS10 6.1]%[ma:1e9bcc0f]
% nguon: Đề tham khảo Tuyển sinh 10 TP.HCM 2025-2026 – Quận 6 – Đề 1 – Bài 3
Cách Tết Trung thu 2 tháng, bà Hoa nhập 100 cái bánh dẻo hạt sen và 150 cái bánh chay mè đen hết $10\,230\,000$ đồng. Thấy bán được, bà Hoa nhập thêm 400 cái bánh dẻo hạt sen và 200 cái bánh chay mè đen (cùng giá nhập đợt đầu) hết $22\,920\,000$ đồng.
\begin{enumerate}[a)]
    \item Hỏi giá nhập của mỗi chiếc bánh dẻo hạt sen và bánh chay mè đen là bao nhiêu?
    \item Trước Tết Trung thu 3 ngày, bà Hoa đã bán hết bánh chay mè đen nhưng vẫn còn 100 cái bánh dẻo hạt sen. Để nhanh chóng bán hết số bánh đã nhập về, bà Hoa thực hiện chương trình khuyến mại như sau: mua cái bánh thứ nhất bằng giá niêm yết, cái bánh thứ hai giảm $30\%$ so với giá niêm yết, cái bánh thứ ba trở đi giảm $60\%$ so với giá niêm yết. Bạn Mai đến cửa hàng bà Hoa mua 5 cái bánh dẻo hạt sen phải trả $145\,000$ đồng. Hỏi giá niêm yết của mỗi cái bánh dẻo hạt sen là bao nhiêu?
\end{enumerate}
\loigiai{
\begin{enumerate}[a)]
    \item Gọi $x$ (đồng) là giá nhập của mỗi chiếc bánh dẻo hạt sen và $y$ (đồng) là giá nhập của mỗi chiếc bánh chay mè đen ($x, y > 0$).\\
    Theo đề bài, chi phí nhập bánh đợt đầu là:
    $100x + 150y = 10\,230\,000$.\\
    Chi phí nhập bánh đợt hai là:
    $400x + 200y = 22\,920\,000$.\\
    Ta có hệ phương trình:
    $\left\{\begin{aligned} &100x + 150y = 10\,230\,000 \\ &400x + 200y = 22\,920\,000 \end{aligned}\right. \Leftrightarrow \left\{\begin{aligned} &400x + 600y = 40\,920\,000 \\ &400x + 200y = 22\,920\,000 \end{aligned}\right. \Leftrightarrow \left\{\begin{aligned} &400y = 18\,000\,000 \\ &100x + 150y = 10\,230\,000 \end{aligned}\right. \Leftrightarrow \left\{\begin{aligned} &y = 45\,000 \\ &x = 34\,800. \end{aligned}\right.$\\
    Vậy giá nhập của mỗi chiếc bánh dẻo hạt sen là $34\,800$ đồng và bánh chay mè đen là $45\,000$ đồng.
    \item Gọi $p$ (đồng) là giá niêm yết của mỗi chiếc bánh dẻo hạt sen ($p > 0$).\\
    Khi bạn Mai mua 5 chiếc bánh dẻo hạt sen theo chương trình khuyến mại:
    \begin{itemize}
        \item Chiếc bánh thứ nhất có giá bằng $100\%$ giá niêm yết: $p$ (đồng).
        \item Chiếc bánh thứ hai được giảm $30\%$, có giá là: $(1 - 30\%)p = 0{,}7p$ (đồng).
        \item Ba chiếc bánh còn lại (từ chiếc thứ ba trở đi) được giảm $60\%$, mỗi chiếc có giá là: $(1 - 60\%)p = 0{,}4p$ (đồng).\\
        Số tiền trả cho 3 chiếc bánh này là: $3 \cdot 0{,}4p = 1{,}2p$ (đồng).
    \end{itemize}
    Tổng số tiền bạn Mai phải trả cho 5 chiếc bánh là:
    $p + 0{,}7p + 1{,}2p = 2{,}9p$ (đồng).\\
    Theo bài ra, số tiền Mai phải trả là $145\,000$ đồng nên ta có phương trình:
    $2{,}9p = 145\,000 \Leftrightarrow p = \dfrac{145\,000}{2{,}9} = 50\,000$ (đồng).\\
    Vậy giá niêm yết của mỗi chiếc bánh dẻo hạt sen là $50\,000$ đồng.
\end{enumerate}
}
\end{bt}

\begin{bt}%[TS10 6.4]%[ma:3d1ee34c]
% nguon: Đề tham khảo Tuyển sinh 10 TP.HCM 2025-2026 – Quận 6 – Đề 2 – Bài 3
Trên cánh đồng có diện tích $160\text{ ha}$ của một đơn vị sản xuất, người ta dành $60\text{ ha}$ để cấy thí điểm giống lúa mới, còn lại vẫn cấy giống lúa cũ. Khi thu hoạch, đầu tiên người ta gặt $8\text{ ha}$ giống lúa cũ và $7\text{ ha}$ giống lúa mới để đối chứng. Kết quả là $7\text{ ha}$ giống lúa mới cho thu hoạch nhiều hơn $8\text{ ha}$ giống lúa cũ là 2 tấn thóc. Biết rằng tổng số thóc (cả hai giống) thu hoạch cả vụ trên $160\text{ ha}$ là 860 tấn. Hỏi năng suất của mỗi giống lúa trên $1\text{ ha}$ là bao nhiêu tấn thóc?
\loigiai{
Diện tích cấy giống lúa cũ là:
$160 - 60 = 100\text{ (ha)}$.\\
Gọi $x$ (tấn/ha) là năng suất của giống lúa cũ và $y$ (tấn/ha) là năng suất của giống lúa mới ($x, y > 0$).\\
Theo đề bài, thu hoạch $7\text{ ha}$ giống lúa mới nhiều hơn $8\text{ ha}$ giống lúa cũ là 2 tấn thóc, ta có phương trình:
$7y - 8x = 2 \Leftrightarrow -8x + 7y = 2$.\\
Tổng sản lượng thóc thu hoạch trên cả cánh đồng $160\text{ ha}$ là 860 tấn, ta có phương trình:
$100x + 60y = 860 \Leftrightarrow 5x + 3y = 43$.\\
Ta có hệ phương trình:
$\left\{\begin{aligned} &-8x + 7y = 2 \\ &5x + 3y = 43 \end{aligned}\right. \Leftrightarrow \left\{\begin{aligned} &-24x + 21y = 6 \\ &35x + 21y = 301 \end{aligned}\right. \Leftrightarrow \left\{\begin{aligned} &59x = 295 \\ &3y = 43 - 5x \end{aligned}\right. \Leftrightarrow \left\{\begin{aligned} &x = 5 \\ &y = 6. \end{aligned}\right.$\\
Các giá trị $x = 5$ và $y = 6$ đều thỏa mãn điều kiện bài toán.\\
Vậy năng suất của giống lúa cũ là 5 tấn/ha và giống lúa mới là 6 tấn/ha.
}
\end{bt}

\begin{bt}%[TS10 6.7]%[ma:c9a01686]
% nguon: Đề tham khảo Tuyển sinh 10 TP.HCM 2025-2026 – Quận 7 – Đề 1 – Bài 6
Thép không gỉ Ferritic là họ thép hợp kim có chứa từ $12\%$ đến $27\%$ crôm. Một nhà máy luyện thép hiện có sẵn một lượng hợp kim thép chứa $10\%$ crôm và một lượng hợp kim thép chứa $30\%$ crôm (các nguyên liệu không bị hao hụt trong quá trình luyện).
\begin{enumerate}[a)]
    \item Tính khối lượng hợp kim thép mỗi loại từ hai loại thép trên dùng để luyện được $500$ tấn thép chứa $16\%$ crôm.
    \item Nhà máy dự định luyện ra loại thép không gỉ Ferritic từ $100$ tấn thép chứa $10\%$ crôm và $x$ tấn thép chứa $30\%$ crôm. Hỏi $x$ nằm trong khoảng nào?
\end{enumerate}
\loigiai{
\begin{enumerate}[a)]
    \item Gọi $x, y$ lần lượt là khối lượng hợp kim thép chứa $10\%$ crôm và chứa $30\%$ crôm cần dùng ($x, y > 0$, đơn vị tấn).\\
    Tổng khối lượng thép cần luyện là $500$ tấn, ta có phương trình:
    $x + y = 500$.\\
    Khối lượng crôm nguyên chất trong $500$ tấn thép chứa $16\%$ crôm là:
    $500 \cdot 16\% = 80\text{ (tấn)}$.\\
    Theo bài ra ta có phương trình:
    $10\% x + 30\% y = 80 \Leftrightarrow 0{,}1x + 0{,}3y = 80 \Leftrightarrow x + 3y = 800$.\\
    Ta có hệ phương trình:
    $\left\{\begin{aligned} &x + y = 500 \\ &x + 3y = 800 \end{aligned}\right. \Leftrightarrow \left\{\begin{aligned} &2y = 300 \\ &x = 500 - y \end{aligned}\right. \Leftrightarrow \left\{\begin{aligned} &y = 150 \\ &x = 350. \end{aligned}\right.$\\
    Vậy cần dùng $350$ tấn thép chứa $10\%$ crôm và $150$ tấn thép chứa $30\%$ crôm.
    \item Khi trộn $100$ tấn thép chứa $10\%$ crôm với $x$ tấn thép chứa $30\%$ crôm ($x > 0$):\\
    Tổng khối lượng hợp kim thu được là: $100 + x$ (tấn).\\
    Tổng khối lượng crôm nguyên chất là:
    $100 \cdot 10\% + x \cdot 30\% = 10 + 0{,}3x$ (tấn).\\
    Tỉ lệ phần trăm crôm trong hợp kim mới là:
    $\dfrac{10 + 0{,}3x}{100 + x}$.\\
    Để hợp kim mới là thép không gỉ Ferritic (chứa từ $12\%$ đến $27\%$ crôm), ta có bất phương trình kép:
    $12\% \le \dfrac{10 + 0{,}3x}{100 + x} \le 27\% \Leftrightarrow 0{,}12 \le \dfrac{10 + 0{,}3x}{100 + x} \le 0{,}27$.\\
    Vì $100 + x > 0$ nên:
    $\left\{\begin{aligned} &0{,}12(100 + x) \le 10 + 0{,}3x \\ &10 + 0{,}3x \le 0{,}27(100 + x) \end{aligned}\right. \Leftrightarrow \left\{\begin{aligned} &12 + 0{,}12x \le 10 + 0{,}3x \\ &10 + 0{,}3x \le 27 + 0{,}27x \end{aligned}\right. \Leftrightarrow \left\{\begin{aligned} &0{,}18x \ge 2 \\ &0{,}03x \le 17 \end{aligned}\right. \Leftrightarrow \left\{\begin{aligned} &x \ge \dfrac{100}{9} \approx 11{,}11 \\ &x \le \dfrac{1700}{3} \approx 566{,}67. \end{aligned}\right.$\\
    Vậy giá trị của $x$ nằm trong khoảng từ $\dfrac{100}{9}$ tấn đến $\dfrac{1700}{3}$ tấn (xấp xỉ từ $11{,}11$ tấn đến $566{,}67$ tấn).
\end{enumerate}
}
\end{bt}

\begin{bt}%[TS10 6.2]%[ma:22050e7c]
% nguon: Đề tham khảo Tuyển sinh 10 TP.HCM 2025-2026 – Quận 7 – Đề 2 – Bài 3
Công ty A sản xuất những chiếc xe lăn cho người khuyết tật với số vốn đầu tư ban đầu là $500$ triệu đồng. Chi phí sản xuất ra một chiếc xe lăn là $2\,500\,000$ đồng và giá bán ra mỗi chiếc là $3\,000\,000$ đồng.
\begin{enumerate}[a)]
    \item Viết hàm số biểu diễn tổng số tiền đã đầu tư đến khi sản xuất ra được $x$ chiếc xe lăn và hàm số biểu diễn số tiền thu được khi bán ra $x$ chiếc xe lăn.
    \item Hỏi công ty cần phải bán bao nhiêu chiếc xe lăn mới có thể thu hồi được số vốn sản xuất ban đầu?
\end{enumerate}
\loigiai{
\begin{enumerate}[a)]
    \item Gọi $x$ là số chiếc xe lăn được sản xuất và bán ra ($x \in \mathbb{N}^*$).
    \begin{itemize}
        \item Tổng số tiền đã đầu tư khi sản xuất ra $x$ chiếc xe lăn gồm số vốn ban đầu và chi phí sản xuất:
        $C(x) = 500\,000\,000 + 2\,500\,000x\text{ (đồng)}$.
        \item Số tiền thu được khi bán ra $x$ chiếc xe lăn là:
        $R(x) = 3\,000\,000x\text{ (đồng)}$.
    \end{itemize}
    \item Để thu hồi được vốn đầu tư ban đầu thì số tiền thu được phải lớn hơn hoặc bằng tổng số tiền đã đầu tư:
    $R(x) \ge C(x) \Leftrightarrow 3\,000\,000x \ge 500\,000\,000 + 2\,500\,000x \Leftrightarrow 500\,000x \ge 500\,000\,000 \Leftrightarrow x \ge 1\,000$.\\
    Vậy công ty cần phải bán được ít nhất $1\,000$ chiếc xe lăn để thu hồi được vốn ban đầu.
\end{enumerate}
}
\end{bt}

\begin{bt}%[TS10 6.1]%[ma:0b4bd70d]
% nguon: Đề tham khảo Tuyển sinh 10 TP.HCM 2025-2026 – Quận 7 – Đề 2 – Bài 5
Trong tháng 11, hai cửa hàng của một thương hiệu thời trang bán được tất cả $1\,100$ sản phẩm. Sang tháng 12, cửa hàng thứ nhất bán vượt mức $15\%$, cửa hàng thứ hai bán vượt mức $20\%$ so với tháng 11, do đó trong tháng 12 cả hai cửa hàng bán được $1\,295$ sản phẩm. Hỏi trong tháng 12 mỗi cửa hàng bán được bao nhiêu sản phẩm?
\loigiai{
Gọi $x, y$ lần lượt là số sản phẩm cửa hàng thứ nhất và cửa hàng thứ hai bán được trong tháng 11 ($x, y \in \mathbb{N}^*$, $x, y < 1\,100$).\\
Trong tháng 11 cả hai cửa hàng bán được $1\,100$ sản phẩm, ta có phương trình:
$x + y = 1\,100$.\\
Số sản phẩm bán thêm trong tháng 12 là:
$1\,295 - 1\,100 = 195\text{ (sản phẩm)}$.\\
Vì cửa hàng thứ nhất vượt mức $15\%$ và cửa hàng thứ hai vượt mức $20\%$ nên ta có phương trình:
$15\% x + 20\% y = 195 \Leftrightarrow 0{,}15x + 0{,}2y = 195 \Leftrightarrow 3x + 4y = 3\,900$.\\
Ta có hệ phương trình:
$\left\{\begin{aligned} &x + y = 1\,100 \\ &3x + 4y = 3\,900 \end{aligned}\right. \Leftrightarrow \left\{\begin{aligned} &4x + 4y = 4\,400 \\ &3x + 4y = 3\,900 \end{aligned}\right. \Leftrightarrow \left\{\begin{aligned} &x = 500 \\ &y = 1\,100 - 500 = 600. \end{aligned}\right.$\\
Các giá trị tìm được đều thỏa mãn điều kiện bài toán.\\
Trong tháng 12:
\begin{itemize}
    \item Cửa hàng thứ nhất bán được: $500 \cdot (1 + 15\%) = 500 \cdot 1{,}15 = 575$ (sản phẩm).
    \item Cửa hàng thứ hai bán được: $600 \cdot (1 + 20\%) = 600 \cdot 1{,}2 = 720$ (sản phẩm).
\end{itemize}
Vậy trong tháng 12, cửa hàng thứ nhất bán được $575$ sản phẩm và cửa hàng thứ hai bán được $720$ sản phẩm.
}
\end{bt}

\begin{bt}%[TS10 6.5]%[ma:f2dd12b9]
% nguon: Đề tham khảo Tuyển sinh 10 TP.HCM 2025-2026 – Quận 8 – Đề 1 – Câu 6
Có ba thùng dầu đựng tổng cộng 123 lít dầu. Nếu đổ từ thùng thứ nhất sang thùng thứ hai 5 lít, rồi đổ từ thùng thứ hai sang thùng thứ ba 7 lít, tiếp tục đổ từ thùng thứ ba sang thùng thứ nhất 9 lít thì số dầu ở thùng thứ nhất sẽ ít hơn số dầu ở thùng thứ hai là 4 lít và bằng $\dfrac{2}{3}$ số dầu ở thùng thứ ba. Tính số lít dầu ở mỗi thùng lúc đầu.
\loigiai{
Gọi số lít dầu ban đầu ở thùng thứ nhất, thùng thứ hai và thùng thứ ba lần lượt là $x, y, z$ ($x, y, z > 0$ và $x + y + z = 123$).\\
Sau khi thực hiện các lần chuyển dầu:
\begin{itemize}
    \item Thùng thứ nhất bớt 5 lít và nhận 9 lít, số dầu còn lại là: $x - 5 + 9 = x + 4$ (lít).
    \item Thùng thứ hai nhận 5 lít và bớt 7 lít, số dầu còn lại là: $y + 5 - 7 = y - 2$ (lít).
    \item Thùng thứ ba nhận 7 lít và bớt 9 lít, số dầu còn lại là: $z + 7 - 9 = z - 2$ (lít).
\end{itemize}
Theo đề bài:
Số dầu thùng thứ nhất ít hơn thùng thứ hai 4 lít, ta có:
$(y - 2) - (x + 4) = 4 \Leftrightarrow y - x - 6 = 4 \Leftrightarrow y = x + 10$.\\
Số dầu thùng thứ nhất bằng $\dfrac{2}{3}$ số dầu thùng thứ ba, ta có:
$x + 4 = \dfrac{2}{3}(z - 2) \Leftrightarrow 3(x + 4) = 2z - 4 \Leftrightarrow 2z = 3x + 16 \Leftrightarrow z = 1{,}5x + 8$.\\
Thay các biểu thức của $y$ và $z$ theo $x$ vào phương trình tổng số dầu:
$x + (x + 10) + (1{,}5x + 8) = 123 \Leftrightarrow 3{,}5x + 18 = 123 \Leftrightarrow 3{,}5x = 105 \Leftrightarrow x = 30$.\\
Từ đó suy ra:
$y = 30 + 10 = 40\text{ (lít)}$,\\
$z = 1{,}5 \cdot 30 + 8 = 53\text{ (lít)}$.\\
Các giá trị tìm được đều thỏa mãn bài toán.\\
Vậy ban đầu thùng thứ nhất có 30 lít dầu, thùng thứ hai có 40 lít dầu và thùng thứ ba có 53 lít dầu.
}
\end{bt}

\begin{bt}%[TS10 6.7]%[ma:2f0421bc]
% nguon: Đề tham khảo Tuyển sinh 10 TP.HCM 2025-2026 – Quận 8 – Đề 2 – Bài 6
Đồng thau là hợp kim có chứa đồng và kẽm. Trong đồng thau màu đỏ có khối lượng kẽm chiếm $20\%$, trong đồng thau màu vàng nhạt có khối lượng kẽm chiếm $42\%$. Từ một ống đồng thau màu đỏ có khối lượng $500\text{ g}$, người ta cần cho thêm bao nhiêu gam kẽm để tạo ra một ống đồng thau màu vàng nhạt? (Làm tròn kết quả đến hàng đơn vị).
\loigiai{
Khối lượng kẽm có sẵn trong $500\text{ g}$ đồng thau màu đỏ là:
$500 \cdot 20\% = 100\text{ (g)}$.\\
Gọi $x$ (g) là khối lượng kẽm cần cho thêm vào ($x > 0$).\\
Khối lượng của hợp kim đồng thau mới sau khi thêm kẽm là:
$500 + x\text{ (g)}$.\\
Khối lượng kẽm trong hợp kim mới là:
$100 + x\text{ (g)}$.\\
Vì hợp kim mới là đồng thau màu vàng nhạt chứa $42\%$ kẽm nên ta có phương trình:
$\dfrac{100 + x}{500 + x} = 42\% \Leftrightarrow \dfrac{100 + x}{500 + x} = 0{,}42 \Leftrightarrow 100 + x = 0{,}42(500 + x) \Leftrightarrow 100 + x = 210 + 0{,}42x \Leftrightarrow 0{,}58x = 110 \Leftrightarrow x = \dfrac{110}{0{,}58} \approx 189{,}66\text{ (g)}$.\\
Làm tròn kết quả đến hàng đơn vị, ta được $x \approx 190\text{ g}$.\\
Vậy cần cho thêm khoảng $190\text{ g}$ kẽm.
}
\end{bt}

\begin{bt}%[TS10 6.1]%[ma:a9906b1f]
% nguon: Đề tham khảo Tuyển sinh 10 TP.HCM 2025-2026 – Quận 10 – Đề 1 – Bài 3
Một cửa hàng bán sỉ nước giải khát Pepsi với giá niêm yết $182\,000$ đồng/thùng. Mỗi tháng cửa hàng xuất bán trung bình được $550$ thùng. Để mở rộng kinh doanh, cửa hàng đã giảm giá $x\%$ cho mỗi thùng Pepsi và đã bán thêm được $y$ thùng trong một tháng.
\begin{enumerate}[a)]
    \item Viết biểu thức số tiền thu được trong một tháng sau khi mở rộng kinh doanh theo $x$ và $y$.
    \item Biết trong tháng đầu tiên sau khi mở rộng, cửa hàng giảm giá $5\%$ và thu được số tiền nhiều hơn khi chưa mở rộng là $38{,}22$ triệu đồng. Hãy tính xem cửa hàng đã bán thêm được bao nhiêu thùng Pepsi trong tháng đó.
\end{enumerate}
\loigiai{
\begin{enumerate}[a)]
    \item Giá bán của mỗi thùng Pepsi sau khi giảm giá $x\%$ là:
    $182\,000 \cdot (1 - x\%)\text{ (đồng)}$.\\
    Số lượng thùng Pepsi bán ra trong một tháng sau khi mở rộng là:
    $550 + y\text{ (thùng)}$.\\
    Biểu thức số tiền thu được trong một tháng theo $x$ và $y$ là:
    $T(x; y) = 182\,000(1 - x\%)(550 + y)\text{ (đồng)}$.
    \item Số tiền cửa hàng thu được mỗi tháng khi chưa mở rộng kinh doanh là:
    $182\,000 \cdot 550 = 100\,100\,000\text{ (đồng)} = 100{,}1\text{ (triệu đồng)}$.\\
    Trong tháng đầu tiên sau khi giảm giá $5\%$, số tiền thu được là:
    $100{,}1 + 38{,}22 = 138{,}32\text{ (triệu đồng)} = 138\,320\,000\text{ (đồng)}$.\\
    Giá bán của mỗi thùng Pepsi sau khi giảm giá $5\%$ là:
    $182\,000 \cdot (1 - 5\%) = 182\,000 \cdot 0{,}95 = 172\,900\text{ (đồng)}$.\\
    Tổng số thùng Pepsi bán ra trong tháng này là:
    $\dfrac{138\,320\,000}{172\,900} = 800\text{ (thùng)}$.\\
    Số thùng Pepsi cửa hàng đã bán thêm được là:
    $y = 800 - 550 = 250\text{ (thùng)}$.\\
    Vậy cửa hàng đã bán thêm được 250 thùng Pepsi.
\end{enumerate}
}
\end{bt}

\begin{bt}%[TS10 6.7]%[ma:c3de256e]
% nguon: Đề tham khảo Tuyển sinh 10 TP.HCM 2025-2026 – Quận 10 – Đề 1 – Bài 5
Một công ty phân bón cần sản xuất một loại phân bón chứa $30\%$ potassium từ hai loại nguyên liệu: loại A chứa $24\%$ potassium và loại B chứa $40\%$ potassium.
\begin{enumerate}[a)]
    \item Tính khối lượng mỗi loại nguyên liệu cần sử dụng để thu được hỗn hợp $500\text{ kg}$ chứa $30\%$ potassium.
    \item Nếu công ty thêm $50\text{ kg}$ loại A vào hỗn hợp $500\text{ kg}$ ban đầu, tỉ lệ phần trăm potassium trong hỗn hợp mới là bao nhiêu (làm tròn số đến hàng phần mười)?
\end{enumerate}
\loigiai{
\begin{enumerate}[a)]
    \item Gọi $x, y$ lần lượt là khối lượng nguyên liệu loại A và loại B cần dùng ($x, y > 0$, đơn vị kg).\\
    Tổng khối lượng hỗn hợp cần sản xuất là $500\text{ kg}$, ta có:
    $x + y = 500$.\\
    Khối lượng potassium nguyên chất trong $500\text{ kg}$ hỗn hợp chứa $30\%$ là:
    $500 \cdot 30\% = 150\text{ (kg)}$.\\
    Theo bài ra ta có phương trình:
    $24\% x + 40\% y = 150 \Leftrightarrow 0{,}24x + 0{,}4y = 150$.\\
    Ta có hệ phương trình:
    $\left\{\begin{aligned} &x + y = 500 \\ &0{,}24x + 0{,}4y = 150 \end{aligned}\right. \Leftrightarrow \left\{\begin{aligned} &0{,}4x + 0{,}4y = 200 \\ &0{,}24x + 0{,}4y = 150 \end{aligned}\right. \Leftrightarrow \left\{\begin{aligned} &0{,}16x = 50 \\ &y = 500 - x \end{aligned}\right. \Leftrightarrow \left\{\begin{aligned} &x = 312{,}5 \\ &y = 187{,}5. \end{aligned}\right.$\\
    Vậy cần dùng $312{,}5\text{ kg}$ loại A và $187{,}5\text{ kg}$ loại B.
    \item Khi thêm $50\text{ kg}$ loại A vào hỗn hợp ban đầu:\\
    Tổng khối lượng của hỗn hợp mới là:
    $500 + 50 = 550\text{ (kg)}$.\\
    Khối lượng potassium trong $50\text{ kg}$ loại A thêm vào là:
    $50 \cdot 24\% = 12\text{ (kg)}$.\\
    Tổng khối lượng potassium trong hỗn hợp mới là:
    $150 + 12 = 162\text{ (kg)}$.\\
    Tỉ lệ phần trăm potassium trong hỗn hợp mới là:
    $\dfrac{162}{550} \cdot 100\% \approx 29{,}45\% \approx 29{,}5\%$.\\
    Vậy hàm lượng potassium trong hỗn hợp mới là khoảng $29{,}5\%$.
\end{enumerate}
}
\end{bt}

\begin{bt}%[TS10 6.2]%[ma:108aa5df]
% nguon: Đề tham khảo Tuyển sinh 10 năm học 2025-2026 Sở GD&ĐT TP.HCM – Quận 10 - Đề 2 (Bài 5)
Hai bạn An và Bình cùng muốn mua một bộ truyện tranh nhưng riêng mỗi bạn thì không đủ tiền.
An nói với Bình: ``Nếu cậu cho tớ mượn $\dfrac{3}{5}$ số tiền cậu có thì tớ sẽ đủ tiền mua bộ truyện''.
Bình nói với An: ``Nếu cậu cho tớ mượn một nửa số tiền của cậu có thì tớ vừa đủ tiền mua bộ truyện''.
Hỏi mỗi bạn có bao nhiêu tiền? Biết rằng số tiền bạn Bình có nhiều hơn số tiền bạn An có là $40\,000$ đồng.
\loigiai{
Gọi số tiền bạn An có là $x$ (đồng) và số tiền bạn Bình có là $y$ (đồng), điều kiện $x > 0, y > 0$.
Vì số tiền bạn Bình có nhiều hơn bạn An là $40\,000$ đồng nên ta có phương trình:
$y - x = 40\,000 \iff -x + y = 40\,000$. \quad (1)
Giá tiền của bộ truyện tranh được tính theo lời An là:
$x + \dfrac{3}{5}y$ (đồng).
Giá tiền của bộ truyện tranh được tính theo lời Bình là:
$y + \dfrac{1}{2}x$ (đồng).
Do giá tiền bộ truyện tranh là không đổi nên:
$x + \dfrac{3}{5}y = y + \dfrac{1}{2}x \iff \dfrac{1}{2}x - \dfrac{2}{5}y = 0 \iff 5x - 4y = 0$. \quad (2)
Từ (1) và (2) ta có hệ phương trình:
$\begin{cases} -x + y = 40\,000 \\ 5x - 4y = 0 \end{cases} \iff \begin{cases} -4x + 4y = 160\,000 \\ 5x - 4y = 0 \end{cases} \iff \begin{cases} x = 160\,000 \\ y = 200\,000 \end{cases}$ (thỏa mãn điều kiện).
Vậy bạn An có $160\,000$ đồng và bạn Bình có $200\,000$ đồng.
}
\end{bt}

\begin{bt}%[TS10 6.8]%[ma:ba7588fb]
% nguon: Đề tham khảo Tuyển sinh 10 năm học 2025-2026 Sở GD&ĐT TP.HCM – Quận 11 - Đề 1 (Bài 3)
Nguyễn Thị Ánh Viên -- kình ngư xuất sắc của bơi lội Việt Nam -- sinh ngày 9 tháng 11 năm $\overline{19ab}$ (với $a, b \in \{0; 1; \dots; 9\}$).
\begin{enumerate}[a)]
    \item Hãy viết biểu thức biểu thị số tuổi của Ánh Viên vào đúng ngày sinh nhật của cô trong năm 2021 theo hai chữ số $a$ và $b$.
    \item Em hãy xác định chính xác năm sinh của Ánh Viên. Biết rằng đến đúng ngày sinh nhật của mình trong năm 2021, cô nhận thấy số tuổi của mình bằng tổng các chữ số trong năm sinh của cô.
\end{enumerate}
\loigiai{
\begin{enumerate}[a)]
    \item Năm sinh của Ánh Viên là $\overline{19ab} = 1900 + 10a + b$.
    Vào đúng ngày sinh nhật 9 tháng 11 năm 2021, số tuổi của cô được biểu thị bởi:
    $T = 2021 - (1900 + 10a + b) = 121 - 10a - b$.
    \item Tổng các chữ số trong năm sinh của Ánh Viên là:
    $S = 1 + 9 + a + b = 10 + a + b$.
    Theo giả thiết, đến ngày sinh nhật năm 2021, số tuổi bằng tổng các chữ số của năm sinh, ta có phương trình:
    $121 - 10a - b = 10 + a + b \iff 11a + 2b = 111$.
    Vì $b \le 9$ nên $2b \le 18 \implies 11a = 111 - 2b \ge 111 - 18 = 93 \implies a \ge \dfrac{93}{11} \approx 8{,}45$.
    Do $a$ là chữ số nên $a = 9$.
    Với $a = 9$, ta có:
    $11 \cdot 9 + 2b = 111 \iff 99 + 2b = 111 \iff 2b = 12 \iff b = 6$ (thỏa mãn $b$ là chữ số).
    Vậy năm sinh của Ánh Viên là năm 1996.
\end{enumerate}
}
\end{bt}

\begin{bt}%[TS10 6.1]%[ma:116d7bb6]
% nguon: Đề tham khảo Tuyển sinh 10 năm học 2025-2026 Sở GD&ĐT TP.HCM – Quận 11 - Đề 1 (Bài 5)
Một trường học tổ chức cho $140$ người gồm giáo viên và học sinh tham gia chuyến tham quan du lịch sinh thái. Giá vé vào cổng niêm yết là $100\,000$ đồng mỗi người. Do được ưu đãi, giá vé vào cổng dành cho học sinh được giảm $10\%$ so với giá niêm yết, do đó tổng số tiền vé nhà trường đã thanh toán là $12\,720\,000$ đồng. Hỏi có bao nhiêu giáo viên và bao nhiêu học sinh đã tham gia chuyến đi?
\loigiai{
Gọi số giáo viên tham gia chuyến đi là $x$ người và số học sinh tham gia là $y$ người ($x, y \in \mathbb{N}^*$, $x, y < 140$).
Vì tổng số người tham gia là $140$ người nên ta có phương trình:
$x + y = 140$. \quad (1)
Giá vé vào cổng của mỗi giáo viên là $100\,000$ đồng.
Giá vé vào cổng của mỗi học sinh sau khi được giảm $10\%$ là:
$100\,000 \cdot (1 - 0{,}10) = 90\,000$ đồng.
Tổng số tiền vé nhà trường đã trả là $12\,720\,000$ đồng nên ta có phương trình:
$100\,000x + 90\,000y = 12\,720\,000 \iff 10x + 9y = 1272$. \quad (2)
Từ (1) và (2) ta có hệ phương trình:
$\begin{cases} x + y = 140 \\ 10x + 9y = 1272 \end{cases} \iff \begin{cases} 9x + 9y = 1260 \\ 10x + 9y = 1272 \end{cases} \iff \begin{cases} x = 12 \\ y = 140 - 12 = 128 \end{cases}$ (thỏa mãn điều kiện).
Vậy có $12$ giáo viên và $128$ học sinh tham gia chuyến đi.
}
\end{bt}

\begin{bt}%[TS10 6.7]%[ma:9e109579]
% nguon: Đề tham khảo Tuyển sinh 10 năm học 2025-2026 Sở GD&ĐT TP.HCM – Quận 11 - Đề 2 (Bài 5)
Theo nghiên cứu dinh dưỡng, lượng protein chứa trong thịt bò và thịt gà lần lượt là $26\%$ và $31\%$; lượng lipit chứa trong thịt bò và thịt gà lần lượt là $15\%$ và $14\%$. Lượng protein và lipit tiêu chuẩn cần thiết cho người trưởng thành mỗi ngày lần lượt là $0{,}8\text{ g/kg}$ và $0{,}9\text{ g/kg}$ trọng lượng cơ thể. Bác Dũng có trọng lượng $70\text{ kg}$. Hỏi bác Dũng cần bổ sung bao nhiêu gam protein và lipit mỗi ngày theo trọng lượng cơ thể? Nếu chỉ dùng thịt bò để bổ sung đủ lượng protein thì cần bao nhiêu gam thịt bò (làm tròn đến hàng đơn vị)?
\loigiai{
Lượng protein cần thiết cho bác Dũng mỗi ngày là:
$m_{\text{protein}} = 70 \cdot 0{,}8 = 56\text{ g}$.
Lượng lipit cần thiết cho bác Dũng mỗi ngày là:
$m_{\text{lipit}} = 70 \cdot 0{,}9 = 63\text{ g}$.
Nếu bác Dũng chỉ ăn thịt bò để bổ sung đủ $56\text{ g}$ protein thì lượng thịt bò cần dùng là:
$m_{\text{bò}} = \dfrac{56}{26\%} = \dfrac{56}{0{,}26} \approx 215\text{ g}$.
}
\end{bt}

\begin{bt}%[TS10 6.5]%[ma:f623ee87]
% nguon: Đề tham khảo Tuyển sinh 10 năm học 2025-2026 Sở GD&ĐT TP.HCM – Quận 12 - Đề 2 (Bài 5)
Hai trường trung học cơ sở $A$ và $B$ có tất cả $250$ học sinh tham gia kỳ thi tuyển sinh vào lớp 10 trung học phổ thông công lập. Kết quả kỳ thi cho thấy: có $\dfrac{2}{3}$ số học sinh dự thi của trường $A$ trúng tuyển và có $\dfrac{3}{5}$ số học sinh dự thi của trường $B$ trúng tuyển. Biết rằng số học sinh trúng tuyển của trường $A$ nhiều hơn số học sinh trúng tuyển của trường $B$ là $2$ học sinh. Hãy tính số học sinh tham gia dự thi của mỗi trường $A$ và $B$.
\loigiai{
Gọi số học sinh tham gia dự thi của trường $A$ là $x$ (học sinh) và của trường $B$ là $y$ (học sinh), điều kiện $x, y \in \mathbb{N}^*$ và $x, y < 250$.
Vì tổng số học sinh dự thi của cả hai trường là $250$ học sinh nên ta có phương trình:
$x + y = 250$. \quad (1)
Số học sinh trúng tuyển của trường $A$ là $\dfrac{2}{3}x$ (học sinh).
Số học sinh trúng tuyển của trường $B$ là $\dfrac{3}{5}y$ (học sinh).
Vì số học sinh trúng tuyển của trường $A$ nhiều hơn trường $B$ là $2$ học sinh nên ta có phương trình:
$\dfrac{2}{3}x - \dfrac{3}{5}y = 2 \iff 10x - 9y = 30$. \quad (2)
Từ (1) và (2) ta có hệ phương trình:
$\begin{cases} x + y = 250 \\ 10x - 9y = 30 \end{cases} \iff \begin{cases} 9x + 9y = 2250 \\ 10x - 9y = 30 \end{cases} \iff \begin{cases} 19x = 2280 \\ y = 250 - x \end{cases} \iff \begin{cases} x = 120 \\ y = 130 \end{cases}$ (thỏa mãn điều kiện).
Vậy trường $A$ có $120$ học sinh và trường $B$ có $130$ học sinh tham gia dự thi.
}
\end{bt}

\begin{bt}%[TS10 6.1]%[ma:572ce528]
% nguon: Đề tham khảo Tuyển sinh 10 năm học 2025-2026 Sở GD&ĐT TP.HCM – TP. Thủ Đức - Đề 1 (Bài 6)
Có hai loại quặng sắt: quặng loại $A$ chứa $60\%$ sắt và quặng loại $B$ chứa $50\%$ sắt. Người ta đem trộn một lượng quặng loại $A$ với một lượng quặng loại $B$ thì thu được một hỗn hợp quặng chứa $\dfrac{8}{15}$ hàm lượng sắt. Nếu lấy tăng thêm $10$ tấn quặng loại $A$ và lấy giảm đi $10$ tấn quặng loại $B$ so với ban đầu thì thu được hỗn hợp quặng mới chứa $\dfrac{17}{30}$ hàm lượng sắt. Hãy tính khối lượng mỗi loại quặng đã đem trộn lúc ban đầu.
\loigiai{
Gọi khối lượng quặng loại $A$ và loại $B$ đem trộn lúc đầu lần lượt là $x$ và $y$ (tấn), điều kiện $x > 0, y > 10$.
Tổng khối lượng quặng đem trộn lúc đầu là $x + y$ (tấn).
Lượng sắt nguyên chất có trong hỗn hợp lúc đầu là:
$0{,}6x + 0{,}5y$ (tấn).
Vì hỗn hợp chứa $\dfrac{8}{15}$ sắt nên ta có phương trình:
$\dfrac{0{,}6x + 0{,}5y}{x + y} = \dfrac{8}{15} \iff 15(0{,}6x + 0{,}5y) = 8(x + y) \iff 9x + 7{,}5y = 8x + 8y \iff x - 0{,}5y = 0 \iff 2x - y = 0$. \quad (1)
Khi tăng $10$ tấn quặng $A$ và giảm $10$ tấn quặng $B$, khối lượng quặng $A$ là $x + 10$ (tấn), quặng $B$ là $y - 10$ (tấn).
Tổng khối lượng quặng hỗn hợp không đổi: $(x + 10) + (y - 10) = x + y$ (tấn).
Lượng sắt trong hỗn hợp lúc sau là:
$0{,}6(x + 10) + 0{,}5(y - 10) = 0{,}6x + 0{,}5y + 1$ (tấn).
Vì hỗn hợp mới chứa $\dfrac{17}{30}$ sắt nên ta có phương trình:
$\dfrac{0{,}6x + 0{,}5y + 1}{x + y} = \dfrac{17}{30} \iff 30(0{,}6x + 0{,}5y + 1) = 17(x + y) \iff 18x + 15y + 30 = 17x + 17y \iff x - 2y = -30$. \quad (2)
Từ (1) và (2) ta có hệ phương trình:
$\begin{cases} 2x - y = 0 \\ x - 2y = -30 \end{cases} \iff \begin{cases} y = 2x \\ x - 2(2x) = -30 \end{cases} \iff \begin{cases} -3x = -30 \\ y = 2x \end{cases} \iff \begin{cases} x = 10 \\ y = 20 \end{cases}$ (thỏa mãn điều kiện).
Vậy khối lượng quặng đem trộn lúc ban đầu là $10$ tấn quặng loại $A$ và $20$ tấn quặng loại $B$.
}
\end{bt}

\begin{bt}%[TS10 6.8]%[ma:cb43fa21]
% nguon: Đề tham khảo Tuyển sinh 10 năm học 2025-2026 Sở GD&ĐT TP.HCM – TP. Thủ Đức - Đề 2 (Bài 6)
Trong một đợt thi khảo sát học tập của lớp 9A, bạn Lan tham gia thi $5$ môn học: Ngữ văn, Toán, Tiếng Anh, Khoa học tự nhiên, Lịch sử và Địa lí. Lan đã hoàn thành $4$ môn với kết quả: Ngữ văn đạt $8$ điểm, Tiếng Anh đạt $7$ điểm, Khoa học tự nhiên đạt $9$ điểm, Lịch sử và Địa lí đạt $7$ điểm. Nhà trường quy định điểm môn Ngữ văn và môn Toán được tính hệ số $2$, ba môn còn lại tính hệ số $1$. Điểm số mỗi môn là các số nguyên từ $1$ đến $10$. Hỏi để đạt được phần thưởng dành cho học sinh có điểm trung bình từ $8{,}0$ điểm trở lên thì Lan phải đạt điểm thi môn Toán ít nhất là bao nhiêu điểm?
\loigiai{
Tổng các hệ số của $5$ môn thi là:
$2 + 2 + 1 + 1 + 1 = 7$.
Gọi điểm thi môn Toán của Lan là $x$ ($x \in \mathbb{N}^*$, $1 \le x \le 10$).
Tổng điểm có tính hệ số của $5$ môn thi là:
$T = 2 \cdot 8 + 2 \cdot x + 7 + 9 + 7 = 2x + 39$.
Điểm trung bình các môn thi của Lan là:
$\overline{X} = \dfrac{2x + 39}{7}$.
Để Lan nhận được phần thưởng, điểm trung bình phải đạt từ $8{,}0$ điểm trở lên, tức là:
$\dfrac{2x + 39}{7} \ge 8 \iff 2x + 39 \ge 56 \iff 2x \ge 17 \iff x \ge 8{,}5$.
Vì điểm thi môn Toán $x$ là một số nguyên dương không vượt quá $10$ nên giá trị nhỏ nhất của $x$ là $9$.
Vậy bạn Lan phải đạt điểm thi môn Toán ít nhất là $9$ điểm để được nhận phần thưởng.
}
\end{bt}

\begin{bt}%[TS10 6.2]%[ma:74d932bb]
% nguon: Đề tham khảo Tuyển sinh 10 năm học 2025-2026 Sở GD&ĐT TP.HCM – TP. Thủ Đức - Đề 3 (Bài 6)
Lớp 9A đăng ký tham gia hoạt động lao động vệ sinh làm sạch khuôn viên trường học. Với số lượng học sinh đăng ký ban đầu, giáo viên chủ nhiệm dự định chia đều thành $3$ tổ có số học sinh bằng nhau. Sau đó, có thêm $4$ học sinh nữa đăng ký tham gia nên cô giáo đã chia đều toàn bộ số học sinh của lớp thành $4$ tổ. Biết rằng mỗi tổ theo cách chia mới có ít hơn $2$ học sinh so với cách chia dự định ban đầu. Hỏi hiện tại lớp 9A có bao nhiêu học sinh tham gia lao động?
\loigiai{
Gọi số học sinh đăng ký tham gia ban đầu là $x$ (học sinh, điều kiện $x \in \mathbb{N}^*$, $x$ chia hết cho 3).
Số học sinh của mỗi tổ theo dự định ban đầu là:
$\dfrac{x}{3}$ (học sinh).
Sau khi có thêm $4$ học sinh, tổng số học sinh của lớp tham gia là:
$x + 4$ (học sinh).
Số học sinh của mỗi tổ theo cách chia mới là:
$\dfrac{x + 4}{4}$ (học sinh).
Vì mỗi tổ hiện nay ít hơn dự định ban đầu $2$ học sinh nên ta có phương trình:
$\dfrac{x}{3} - \dfrac{x + 4}{4} = 2 \iff 4x - 3(x + 4) = 24 \iff x - 12 = 24 \iff x = 36$ (thỏa mãn điều kiện).
Số học sinh hiện tại của lớp tham gia lao động là:
$36 + 4 = 40$ (học sinh).
Vậy hiện tại lớp 9A có $40$ học sinh tham gia hoạt động vệ sinh trường học.
}
\end{bt}

\begin{bt}%[TS10 6.1]%[ma:f8535bde]
% nguon: Đề tham khảo Tuyển sinh 10 năm học 2025-2026 Sở GD&ĐT TP.HCM – Quận Tân Phú - Đề 1 (Bài 5)
Theo kế hoạch, hai tổ sản xuất được giao nhiệm vụ hoàn thành tổng cộng $700$ sản phẩm trong một thời gian nhất định. Nhờ áp dụng cải tiến kỹ thuật, tổ I đã sản xuất vượt mức $12\%$ so với kế hoạch, trong khi tổ II gặp khó khăn nên sản lượng giảm $5\%$ so với kế hoạch. Tuy nhiên, tính chung cả hai tổ vẫn hoàn thành vượt mức $16$ sản phẩm. Hãy tính số sản phẩm làm vượt mức của tổ I và số sản phẩm bị hụt giảm của tổ II so với chỉ tiêu kế hoạch ban đầu.
\loigiai{
Gọi số sản phẩm được giao theo kế hoạch của tổ I là $x$ và của tổ II là $y$ ($x, y \in \mathbb{N}^*$, $x, y < 700$).
Vì tổng số sản phẩm theo kế hoạch của hai tổ là $700$ nên ta có phương trình:
$x + y = 700$. \quad (1)
Số sản phẩm tổ I làm vượt mức là: $12\% \cdot x = 0{,}12x$ (sản phẩm).
Số sản phẩm tổ II bị giảm là: $5\% \cdot y = 0{,}05y$ (sản phẩm).
Do cả hai tổ hoàn thành vượt mức $16$ sản phẩm so với kế hoạch nên:
$0{,}12x - 0{,}05y = 16$. \quad (2)
Từ (1) và (2) ta có hệ phương trình:
$\begin{cases} x + y = 700 \\ 0{,}12x - 0{,}05y = 16 \end{cases} \iff \begin{cases} 0{,}05x + 0{,}05y = 35 \\ 0{,}12x - 0{,}05y = 16 \end{cases} \iff \begin{cases} 0{,}17x = 51 \\ y = 700 - x \end{cases} \iff \begin{cases} x = 300 \\ y = 400 \end{cases}$ (thỏa mãn điều kiện).
Số sản phẩm tổ I làm vượt mức là: $0{,}12 \cdot 300 = 36$ (sản phẩm).
Số sản phẩm tổ II bị giảm hụt là: $0{,}05 \cdot 400 = 20$ (sản phẩm).
}
\end{bt}

\begin{bt}%[TS10 6.5]%[ma:299475b7]
% nguon: Đề tham khảo Tuyển sinh 10 năm học 2025-2026 Sở GD&ĐT TP.HCM – Quận Tân Phú - Đề 2 (Bài 6)
Trong một kỳ thi thử tốt nghiệp THCS gồm $50$ câu hỏi trắc nghiệm, mỗi câu trả lời đúng được cộng $20$ điểm, còn mỗi câu trả lời sai bị trừ đi $5$ điểm. Bạn Na đã tham gia làm trọn vẹn cả $50$ câu hỏi và đạt được tổng số điểm là $550$ điểm. Hỏi bạn Na đã trả lời đúng bao nhiêu câu và trả lời sai bao nhiêu câu?
\loigiai{
Gọi số câu trả lời đúng của Na là $x$ và số câu trả lời sai là $y$ ($x, y \in \mathbb{N}$, $x, y \le 50$).
Vì Na làm trọn vẹn $50$ câu hỏi nên ta có phương trình:
$x + y = 50$. \quad (1)
Mỗi câu đúng được $+20$ điểm và mỗi câu sai bị $-5$ điểm, tổng số điểm đạt được là $550$ điểm nên:
$20x - 5y = 550 \iff 4x - y = 110$. \quad (2)
Cộng hai phương trình (1) và (2) theo từng vế, ta được:
$5x = 160 \iff x = 32$.
Thay $x = 32$ vào (1), ta được:
$y = 50 - 32 = 18$ (thỏa mãn điều kiện).
Vậy bạn Na đã trả lời đúng $32$ câu và trả lời sai $18$ câu.
}
\end{bt}
